%%%%%%%%%%%%%%%%%%%%%%%%%%%%%%%%%%%%%%%%%%%%%%%%%%%%%%%%%%%%%%%%%
% PHIẾU HỌC TẬP - MÔ HÌNH TOÁN HỌC VỀ LÂY LAN DỊCH BỆNH
% Phiên bản 3.1 - Hoàn chỉnh cho Cuộc thi STEM/STEAM
%%%%%%%%%%%%%%%%%%%%%%%%%%%%%%%%%%%%%%%%%%%%%%%%%%%%%%%%%%%%%%%%%
\documentclass[12pt,a4paper]{article}

\usepackage[utf8]{inputenc}
\usepackage[T1]{fontenc}
\usepackage[vietnamese]{babel}
\usepackage{amsmath,amssymb}
\usepackage{array}
\usepackage{booktabs}
\usepackage{tabularx}
\usepackage{colortbl}
\usepackage{graphicx}
\usepackage{newtxtext, newtxmath}
\usepackage[scaled=0.92]{helvet}
\usepackage{microtype}
\usepackage{tikz}
\usetikzlibrary{shapes,arrows,positioning}

\usepackage[svgnames]{xcolor}
\definecolor{primary}{RGB}{26,95,122}
\definecolor{mathblue}{RGB}{0,51,102}
\definecolor{accent}{RGB}{194,55,40}
\definecolor{success}{RGB}{34,139,34}
\definecolor{TableHeader}{gray}{0.88}
\definecolor{lightgray}{gray}{0.95}

\usepackage[left=1.5cm,right=1.5cm,top=2cm,bottom=2cm,headheight=15pt]{geometry}
\usepackage{fancyhdr}
\pagestyle{fancy}
\fancyhf{}
\renewcommand{\headrule}{\color{primary}\hrule width\textwidth height 1pt}
\fancyhead[L]{\color{primary}\footnotesize\sffamily\textbf{MÔ HÌNH TOÁN HỌC VỀ LÂY LAN DỊCH BỆNH}}
\fancyhead[R]{\color{primary}\footnotesize\sffamily\textbf{TOÁN 11}}
\fancyfoot[C]{\sffamily\thepage}

\usepackage{enumitem}
\usepackage{tcolorbox}
\tcbuselibrary{skins, breakable}
\usepackage[colorlinks=true,linkcolor=primary,urlcolor=accent]{hyperref}

% ============= MÔI TRƯỜNG =============
\newtcolorbox{phieuhoctap}[1]{
    enhanced, breakable,
    colback=white, colframe=primary,
    boxrule=1.5pt, arc=4mm,
    fonttitle=\bfseries\sffamily\Large,
    coltitle=white, colbacktitle=primary,
    title=#1,
    shadow={2mm}{-1mm}{0mm}{black!20}
}

\newtcolorbox{mathbox}{
    enhanced, breakable,
    colback=mathblue!5, colframe=mathblue!80,
    boxrule=1.5pt, arc=3mm,
    fonttitle=\bfseries\sffamily,
    title=\color{white}CÔNG THỨC TOÁN HỌC,
    colbacktitle=mathblue
}

\newtcolorbox{problembox}{
    enhanced, breakable,
    colback=accent!5, colframe=accent!80,
    boxrule=1.5pt, arc=3mm,
    fonttitle=\bfseries\sffamily,
    title=\color{white}BÀI TOÁN DỊCH BỆNH,
    colbacktitle=accent
}

\newtcolorbox{khaosat}[1]{
    enhanced, breakable,
    colback=lightgray, colframe=gray!60,
    boxrule=1pt, arc=3mm,
    fonttitle=\bfseries\sffamily\large,
    coltitle=primary, colbacktitle=lightgray,
    title=#1
}

\newcolumntype{C}{>{\centering\arraybackslash}X}
\newcolumntype{L}{>{\raggedright\arraybackslash}X}
\newcommand{\header}[1]{\cellcolor{TableHeader}\sffamily\bfseries#1}

\setlist[itemize,1]{label=\color{primary}\textbullet, leftmargin=*, nosep}
\setlist[enumerate,1]{label=\color{primary}\arabic*., leftmargin=*, nosep}

\begin{document}

%==============================================================
% KHẢO SÁT TRƯỚC BÀI HỌC
%==============================================================
\begin{khaosat}{PHIẾU KHẢO SÁT TRƯỚC BÀI HỌC}

\noindent\textbf{Họ và tên:} \underline{\hspace{5cm}} \hfill \textbf{Lớp:} \underline{\hspace{2cm}} \hfill \textbf{Ngày:} \underline{\hspace{2cm}}

\vspace{0.5cm}
\textbf{Hướng dẫn:} Đánh dấu ($\checkmark$) vào ô phù hợp với mức độ đồng ý của em.

\vspace{0.3cm}
\begin{tabularx}{\textwidth}{|L|c|c|c|c|}
\hline
\header{Câu hỏi} & \header{5} & \header{4} & \header{3} & \header{1-2} \\
\hline
1. Em hiểu dịch bệnh lây lan như thế nào & $\square$ & $\square$ & $\square$ & $\square$ \\
\hline
2. Em biết hệ số $R_0$ (hệ số lây nhiễm cơ bản) là gì & $\square$ & $\square$ & $\square$ & $\square$ \\
\hline
3. Em hiểu ``miễn dịch cộng đồng'' nghĩa là gì & $\square$ & $\square$ & $\square$ & $\square$ \\
\hline
4. Em biết Hàm số mũ có thể ứng dụng trong y tế & $\square$ & $\square$ & $\square$ & $\square$ \\
\hline
5. Em thấy Toán học giúp dự đoán diễn biến dịch bệnh & $\square$ & $\square$ & $\square$ & $\square$ \\
\hline
\end{tabularx}

\textit{\footnotesize Thang điểm: 5 = Rất đồng ý, 4 = Đồng ý, 3 = Phân vân, 1-2 = Không đồng ý}

\vspace{0.4cm}
\textbf{6. Em mong muốn học được gì từ bài học này?}

\underline{\hspace{0.95\textwidth}}
\end{khaosat}

\vspace{0.8cm}

%==============================================================
% PHIẾU HỌC TẬP SỐ 1
%==============================================================
\begin{phieuhoctap}{PHIẾU HỌC TẬP SỐ 1: XÂY DỰNG CÔNG THỨC TOÁN HỌC}

\noindent\textbf{Họ và tên:} \underline{\hspace{5cm}} \hfill \textbf{Nhóm:} \underline{\hspace{2cm}}

\vspace{0.5cm}
\begin{mathbox}
\textbf{Kiến thức cần nhớ:}

\textbf{1. Hàm số mũ cơ sở $e$:} $y = e^x$ với $e \approx 2.71828$

\textbf{2. Công thức lây lan dịch bệnh:}
$$\boxed{I(t) = I_0 \cdot e^{rt} \quad \text{với} \quad r = \beta - \gamma}$$

\textbf{3. Hệ số lây nhiễm:} $R_0 = \dfrac{\beta}{\gamma}$ \qquad \textbf{4. Ngưỡng miễn dịch:} $H = 1 - \dfrac{1}{R_0}$
\end{mathbox}

\vspace{0.5cm}
\textbf{\color{mathblue} PHẦN A: TÍNH GIÁ TRỊ $e^x$} (Dùng máy tính, làm tròn 2 chữ số)

\begin{center}
\begin{tabularx}{0.9\textwidth}{|C|C|C|C|C|C|}
\hline
$x$ & 0 & 1 & 2 & 3 & 4 \\
\hline
$e^x$ & \underline{\hspace{1.2cm}} & \underline{\hspace{1.2cm}} & \underline{\hspace{1.2cm}} & \underline{\hspace{1.2cm}} & \underline{\hspace{1.2cm}} \\
\hline
\end{tabularx}
\end{center}

\vspace{0.3cm}
\begin{problembox}
\textbf{Tình huống:} Một thành phố xuất hiện \textbf{10 ca nhiễm} bệnh X ban đầu.

Biết: $\beta = 0.3$ (tỷ lệ lây nhiễm/ngày), $\gamma = 0.1$ (tỷ lệ hồi phục/ngày)
\end{problembox}

\vspace{0.3cm}
\textbf{\color{mathblue} PHẦN B: XÁC ĐỊNH THAM SỐ}
\begin{itemize}
\item Số ca ban đầu: $I_0 = \underline{\hspace{2cm}}$
\item Tốc độ lây lan: $r = \beta - \gamma = \underline{\hspace{1.5cm}} - \underline{\hspace{1.5cm}} = \underline{\hspace{2cm}}$
\item Công thức: $I(t) = \underline{\hspace{1.5cm}} \cdot e^{\underline{\hspace{1.5cm}} \cdot t}$
\end{itemize}

\vspace{0.3cm}
\textbf{\color{mathblue} PHẦN C: TÍNH SỐ CA NHIỄM} (Công thức: $I(t) = 10 \cdot e^{0.2t}$)

\begin{center}
\begin{tabularx}{0.95\textwidth}{|C|C|C|C|C|C|}
\hline
\header{Ngày $t$} & 0 & 5 & 10 & 15 & 20 \\
\hline
$0.2t$ & 0 & \underline{\hspace{1cm}} & \underline{\hspace{1cm}} & \underline{\hspace{1cm}} & \underline{\hspace{1cm}} \\
\hline
$e^{0.2t}$ & 1 & \underline{\hspace{1cm}} & \underline{\hspace{1cm}} & \underline{\hspace{1cm}} & \underline{\hspace{1cm}} \\
\hline
$I(t)$ = số ca & 10 & \underline{\hspace{1cm}} & \underline{\hspace{1cm}} & \underline{\hspace{1cm}} & \underline{\hspace{1cm}} \\
\hline
\end{tabularx}
\end{center}

\textbf{Nhận xét:} Số ca tăng rất nhanh vì \underline{\hspace{7cm}}

\vspace{0.3cm}
\textbf{\color{mathblue} PHẦN D: TÍNH $R_0$ VÀ NGƯỠNG MIỄN DỊCH $H$}
\begin{itemize}
\item $R_0 = \dfrac{\beta}{\gamma} = \dfrac{\underline{\hspace{1cm}}}{\underline{\hspace{1cm}}} = \underline{\hspace{2cm}}$
\item Ý nghĩa $R_0$: Mỗi người nhiễm trung bình lây cho \underline{\hspace{2cm}} người
\item Nếu $R_0 > 1$, dịch sẽ \underline{\hspace{4cm}} (bùng phát / tự tắt)
\item $H = 1 - \dfrac{1}{R_0} = 1 - \dfrac{1}{\underline{\hspace{1cm}}} = \underline{\hspace{2cm}} = \underline{\hspace{2cm}}\%$
\item Ý nghĩa $H$: Cần \underline{\hspace{2cm}}\% dân số có miễn dịch để dịch tự tắt
\end{itemize}

\end{phieuhoctap}

\newpage
%==============================================================
% PHIẾU HỌC TẬP SỐ 2
%==============================================================
\begin{phieuhoctap}{PHIẾU HỌC TẬP SỐ 2: MÔ PHỎNG VỚI EPIDEMICSIM}

\noindent\textbf{Họ và tên:} \underline{\hspace{5cm}} \hfill \textbf{Nhóm:} \underline{\hspace{2cm}}

\vspace{0.5cm}
\textbf{\color{primary}\large I. THỰC HÀNH MÔ PHỎNG}

\begin{itemize}
\item \textbf{Địa chỉ:} \href{https://stem-1h8.pages.dev/sim.html}{\texttt{stem-1h8.pages.dev/sim.html}}
\item \textbf{Nhập tham số:} $N = 10{,}000$, $I_0 = 10$, $\beta = 0.3$, $\gamma = 0.1$, Tiêm chủng = 0\%
\end{itemize}

\textbf{\color{accent} Bài 1: So sánh công thức với mô phỏng SIR}

\begin{center}
\begin{tabularx}{0.9\textwidth}{|L|C|C|C|C|C|}
\hline
\header{Ngày $t$} & 0 & 5 & 10 & 15 & 20 \\
\hline
Công thức $I = 10e^{0.2t}$ & 10 & 27 & 74 & 201 & 546 \\
\hline
Từ EpidemicSim & 10 & \underline{\hspace{1cm}} & \underline{\hspace{1cm}} & \underline{\hspace{1cm}} & \underline{\hspace{1cm}} \\
\hline
\end{tabularx}
\end{center}

\textbf{Câu hỏi:}
\begin{enumerate}[label=\alph*)]
\item Giai đoạn đầu (0-10 ngày): Công thức và mô phỏng \underline{\hspace{3cm}} (giống / khác)
\item Giai đoạn sau (15-20 ngày): Kết quả khác nhau vì \underline{\hspace{5cm}}
\item Đỉnh dịch (số ca cao nhất) xảy ra vào ngày \underline{\hspace{2cm}} với \underline{\hspace{3cm}} ca
\item Cuối cùng, tổng số người đã nhiễm (R) là \underline{\hspace{3cm}} người
\end{enumerate}

\vspace{0.5cm}
\textbf{\color{primary}\large II. SO SÁNH KỊCH BẢN CAN THIỆP}

\begin{center}
\begin{tabularx}{0.95\textwidth}{|L|c|c|C|C|C|}
\hline
\header{Kịch bản} & \header{$\beta$} & \header{Tiêm} & \header{Đỉnh (ca)} & \header{Ngày đỉnh} & \header{Tổng R} \\
\hline
KB1: Không can thiệp & 0.3 & 0\% & \underline{\hspace{1.5cm}} & \underline{\hspace{1cm}} & \underline{\hspace{1.5cm}} \\
\hline
KB2: Giãn cách xã hội & 0.15 & 0\% & \underline{\hspace{1.5cm}} & \underline{\hspace{1cm}} & \underline{\hspace{1.5cm}} \\
\hline
KB3: Tiêm chủng 60\% & 0.3 & 60\% & \underline{\hspace{1.5cm}} & \underline{\hspace{1cm}} & \underline{\hspace{1.5cm}} \\
\hline
\end{tabularx}
\end{center}

\textbf{\color{accent} Bài 2: Phân tích kết quả}
\begin{enumerate}[label=\alph*)]
\item Kịch bản hiệu quả nhất là: \underline{\hspace{8cm}}
\item Giãn cách xã hội làm giảm $\beta$ vì \underline{\hspace{6cm}}
\item Với $R_0 = 3$, ngưỡng $H = 66.7\%$. Tiêm 60\% có đủ không? \underline{\hspace{3cm}}
\end{enumerate}

\vspace{0.5cm}
\textbf{\color{primary}\large III. ĐỀ XUẤT CHÍNH SÁCH}

Dựa trên kết quả mô phỏng, nhóm em đề xuất biện pháp phòng dịch:

\underline{\hspace{0.95\textwidth}}

\underline{\hspace{0.95\textwidth}}

\underline{\hspace{0.95\textwidth}}

\end{phieuhoctap}

\newpage
%==============================================================
% KHẢO SÁT SAU BÀI HỌC
%==============================================================
\begin{khaosat}{PHIẾU KHẢO SÁT SAU BÀI HỌC}

\noindent\textbf{Họ và tên:} \underline{\hspace{5cm}} \hfill \textbf{Lớp:} \underline{\hspace{2cm}} \hfill \textbf{Ngày:} \underline{\hspace{2cm}}

\vspace{0.5cm}
\textbf{Hướng dẫn:} Đánh dấu ($\checkmark$) vào ô phù hợp.

\vspace{0.3cm}
\begin{tabularx}{\textwidth}{|L|c|c|c|c|}
\hline
\header{Câu hỏi} & \header{5} & \header{4} & \header{3} & \header{1-2} \\
\hline
1. Em hiểu công thức $I(t) = I_0 \cdot e^{rt}$ & $\square$ & $\square$ & $\square$ & $\square$ \\
\hline
2. Em biết ý nghĩa của $R_0$ và cách tính & $\square$ & $\square$ & $\square$ & $\square$ \\
\hline
3. Em hiểu ngưỡng miễn dịch cộng đồng $H$ & $\square$ & $\square$ & $\square$ & $\square$ \\
\hline
4. Ứng dụng EpidemicSim giúp em hiểu bài hơn & $\square$ & $\square$ & $\square$ & $\square$ \\
\hline
5. Em thấy Toán học có ích trong đời sống & $\square$ & $\square$ & $\square$ & $\square$ \\
\hline
6. Em tự tin giải thích được cách phòng dịch & $\square$ & $\square$ & $\square$ & $\square$ \\
\hline
\end{tabularx}

\textit{\footnotesize Thang điểm: 5 = Rất đồng ý, 4 = Đồng ý, 3 = Phân vân, 1-2 = Không đồng ý}

\vspace{0.4cm}
\textbf{7. Điều em thích nhất về bài học này:}

\underline{\hspace{0.95\textwidth}}

\textbf{8. Điều em muốn cải thiện:}

\underline{\hspace{0.95\textwidth}}
\end{khaosat}

\vspace{0.8cm}

%==============================================================
% RUBRIC ĐÁNH GIÁ
%==============================================================
\begin{center}
{\Large\bfseries\sffamily\color{primary} RUBRIC ĐÁNH GIÁ SẢN PHẨM NHÓM}
\end{center}

\vspace{0.3cm}
\noindent\textbf{Tên nhóm:} \underline{\hspace{4cm}} \hfill \textbf{Người đánh giá:} \underline{\hspace{4cm}}

\vspace{0.5cm}
\begin{tabularx}{\textwidth}{|>{\raggedright\arraybackslash}p{2.2cm}|X|X|X|c|}
\hline
\rowcolor{TableHeader}
\sffamily\bfseries Tiêu chí & \sffamily\bfseries Xuất sắc (4đ) & \sffamily\bfseries Khá (3đ) & \sffamily\bfseries Cần cải thiện (1-2đ) & \sffamily\bfseries Điểm \\
\hline
\textbf{1. Công thức Toán} & Viết đúng $I=I_0 e^{rt}$, tính chính xác, giải thích rõ & Viết đúng, tính đúng & Còn sai công thức & \underline{\hspace{0.8cm}} \\
\hline
\textbf{2. Hiểu $R_0$, $H$} & Giải thích đúng ý nghĩa, tính chính xác, liên hệ thực tế & Hiểu cơ bản, tính đúng & Chưa hiểu ý nghĩa & \underline{\hspace{0.8cm}} \\
\hline
\textbf{3. Mô phỏng} & Thành thạo, so sánh 3 kịch bản, phân tích sâu sắc & Sử dụng tốt, so sánh 2 kịch bản & Cần hướng dẫn & \underline{\hspace{0.8cm}} \\
\hline
\textbf{4. Trình bày} & Mạch lạc, có đồ thị, trả lời phản biện tốt & Rõ ràng, có số liệu & Rời rạc, thiếu minh họa & \underline{\hspace{0.8cm}} \\
\hline
\end{tabularx}

\vspace{0.3cm}
\noindent\textbf{TỔNG ĐIỂM:} \underline{\hspace{1.5cm}} / 16 điểm \hfill \textbf{Xếp loại:} \underline{\hspace{3cm}}

\vspace{0.3cm}
\noindent\textbf{Nhận xét chung:}

\underline{\hspace{0.95\textwidth}}

\underline{\hspace{0.95\textwidth}}

\newpage
%==============================================================
% PHIẾU ĐÁNH GIÁ ĐỒNG CẤP
%==============================================================
\begin{center}
{\Large\bfseries\sffamily\color{primary} PHIẾU ĐÁNH GIÁ ĐỒNG CẤP (Peer Assessment)}
\end{center}

\vspace{0.3cm}
\noindent\textbf{Tên em:} \underline{\hspace{4cm}} \hfill \textbf{Nhóm:} \underline{\hspace{2cm}} \hfill \textbf{Ngày:} \underline{\hspace{2cm}}

\vspace{0.5cm}
\textbf{Hướng dẫn:} Đánh giá các thành viên trong nhóm theo thang điểm 1-5 (1 = Yếu, 5 = Xuất sắc)

\vspace{0.3cm}
\begin{tabularx}{\textwidth}{|L|c|c|c|c|}
\hline
\rowcolor{TableHeader}
\sffamily\bfseries Tên thành viên & \sffamily\bfseries Ý tưởng & \sffamily\bfseries Hợp tác & \sffamily\bfseries Hoàn thành & \sffamily\bfseries Tổng \\
\hline
\underline{\hspace{4cm}} & \underline{\hspace{1cm}} & \underline{\hspace{1cm}} & \underline{\hspace{1cm}} & \underline{\hspace{1cm}} \\
\hline
\underline{\hspace{4cm}} & \underline{\hspace{1cm}} & \underline{\hspace{1cm}} & \underline{\hspace{1cm}} & \underline{\hspace{1cm}} \\
\hline
\underline{\hspace{4cm}} & \underline{\hspace{1cm}} & \underline{\hspace{1cm}} & \underline{\hspace{1cm}} & \underline{\hspace{1cm}} \\
\hline
\underline{\hspace{4cm}} & \underline{\hspace{1cm}} & \underline{\hspace{1cm}} & \underline{\hspace{1cm}} & \underline{\hspace{1cm}} \\
\hline
\end{tabularx}

\vspace{0.3cm}
\textit{\footnotesize Tiêu chí: Ý tưởng = Đóng góp sáng tạo | Hợp tác = Làm việc nhóm | Hoàn thành = Hoàn thành nhiệm vụ đúng hạn}

\vspace{0.5cm}
\textbf{Thành viên đóng góp nhiều nhất:} \underline{\hspace{4cm}}

\textbf{Lý do:} \underline{\hspace{10cm}}

\end{document}
